\chapter{Discussion}\label{ch:discussion}


\section{Performance and Economic Trade-offs}

Across the four datasets evaluated here, the EBM shows no uniform performance
disadvantage relative to the boosted ensembles. On Taiwan its ROC-AUC is effectively
level with XGBoost and CatBoost, and all four paired EBM--ensemble intervals contain
zero.

The Home Credit result is of a different kind. Here the EBM--CatBoost difference in
ROC-AUC lies entirely above zero, so on the second-largest test partition the only
EBM--ensemble contrast that resolves runs in favour of the glass-box model rather than
against it. It is the only such result in ROC-AUC, and it is narrow: the margin is
$0.0016$, the companion contrast against XGBoost includes zero at its lower bound, and
neither EMP interval resolves.

What it establishes is that a resolved direction in favour of the glass-box model
exists. Home Credit and Lending Club resolve in opposite directions, and the remaining
two resolve nothing. This supports the claim that an additive
model with a bounded number of fitted pairwise interactions can reach ensemble-level
discrimination in some tabular credit settings, consistent with prior reports
\citep{changHowInterpretableTrustworthy2021, dessainCostExplainabilityAI2023}. 

German Credit and Lending Club mark different limits to that conclusion. German Credit
provides too little test-sample precision to establish a stable statistical or economic
ordering: no paired interval for any metric excludes zero. Lending Club is the clear counterexample. XGBoost and CatBoost
lead the EBM by about $0.008$ ROC-AUC and by $0.047$--$0.050$ percentage points EMP,
with both EBM--ensemble contrasts resolved in ROC-AUC and in EMP under test-case
resampling.


EMP was included as a second evaluation dimension because ranking quality alone does not
capture the asymmetric consequences of credit decisions. On these datasets it largely
reproduces the statistical evaluation rather than adding an ordering of its own.

The EMP ordering matches the ROC-AUC ordering exactly on Home Credit and Lending
Club, these are the only two datasets in which an EBM--ensemble contrast resolves a direction. It
departs from that ordering only where the performance differences are unresolved to
begin with. Even where a direction does resolve, EMP acts as the more conservative
measure rather than an independent signal: on Home Credit the EBM--CatBoost contrast
excludes zero in ROC-AUC but not in EMP. Nowhere does the profit measure reverse a
direction that discrimination has already resolved.

Its contribution here is therefore not necessarily a second ordering but a second unit. A gap of
$0.008$ ROC-AUC cannot be weighed against anything on its own, but expressed in percentage
points of normalised principal, the same gap becomes a quantity an institution can hold
against the cost of the alternative.

Under the EMP parameterisation, the observed Lending Club difference can be read as a
modelled economic price. The EBM's EMP is $0.047$--$0.050$ percentage points of
normalised principal below that of the two ensembles. This corresponds to approximately
EUR~$0.50$ per EUR~$1{,}000$ of normalised principal, compared with an EMP level of
about EUR~$24.80$ per EUR~$1{,}000$. The difference therefore amounts to approximately
$2\%$ of the modelled EMP level.

Because EMP levels are not directly comparable across datasets, the interpretation
focuses on paired model differences within each dataset. Only Lending Club yields paired
EBM--ensemble EMP intervals that exclude zero. On Home Credit, the point estimates favour
the EBM, but both paired difference intervals include zero and neither contrast is resolved
on Taiwan. Profit-based
evaluation therefore identifies a modelled economic disadvantage for the EBM in only one
of four settings, while its direction remains unresolved in the other three.

Because all configurations are tuned to maximise validation ROC-AUC, this high agreement between the metrics is partly induced by the experimental design. The EMP figures therefore serve as an economic lens applied to standard, discrimination-optimised pipelines, quantifying the economic consequences of models chosen solely for their statistical ranking power.

Logistic regression serves an important comparative role by helping to distinguish
intrinsic interpretability from the restrictions of a linear functional form. Although
it attains the highest threshold-metric point estimates on the small German Credit
dataset, it is consistently less competitive with the non-linear models across the 
larger datasets. This pattern shows that direct inspectability alone is insufficient 
for model choice. An interpretable model must also provide sufficient functional 
flexibility for the empirical setting. The EBM therefore provides the more informative 
test of whether inherent interpretability can exist alongside competitive performance.

\section{Explanation Stability, Agreement and Interpretability}

Under case resampling, top-ten membership is high across all four pipelines and partly
compressed by the small feature vocabularies of German Credit and Taiwan. The more
discriminating rank-interval widths lead to the same substantive conclusion: neither LR
nor the EBM shows a consistent stability advantage over XGBoost and CatBoost.

Cross-model agreement provides a different form of evidence. Top-ten Jaccard indicates
partial overlap between the features prioritised by the pipelines, while the complete
rankings place the EBM close to the two ensembles; logistic regression diverges most
strongly on Taiwan. This agreement indicates similarity in global feature priorities,
but it does not identify a correct or faithful ranking, because attribution allocation
depends on the fitted model, explainer semantics and feature dependence
\citep{aasExplainingIndividualPredictions2021}.

For model choice, the explanation measures therefore add two qualifications. They
provide no stability-based reason to prefer the intrinsically interpretable pipelines,
but they show that the EBM and the ensembles often emphasise similar source features.
The relevant interpretability distinction remains structural: LR and the EBM expose
components of their fitted prediction functions, whereas XGBoost and CatBoost require a
separate post-hoc explanation computation
\citep{rudinStopExplainingBlack2019a}. The case for the EBM consequently rests on this
direct inspectability together with its observed performance, not on superior measured
attribution stability.


\section{Sensitivity to Training Size and Time}

The learning curves show that relative model performance changes with training-set
size under the fixed configurations examined. XGBoost and CatBoost gain more
ROC-AUC as the sample grows. The EBM remains ahead of logistic regression, while
the ensembles close or reverse their initial deficits. The EBM is therefore more
competitive at smaller sample sizes in these experiments, whereas the boosted-tree
ensembles benefit more from additional data. Across-refit top-ten Jaccard overlap
does not increase monotonically, indicating that larger training samples do not
consistently produce more reproducible global feature rankings.

The out-of-time evaluation yields lower ROC-AUC estimates for all Lending Club
pipelines but preserves their random-split ordering: XGBoost and CatBoost lead,
followed by the EBM and logistic regression. The same point-estimate hierarchy
appears for EMP. The chronological evaluation therefore reproduces the central
Lending Club pattern in both discrimination and economic performance. The
EBM--ensemble ROC-AUC gap is nevertheless narrower than under the random split:
the direction of the observed difference recurs, while its magnitude changes.
Taken together, these analyses show that key elements of the model ordering remain across alternative settings: the EBM remains ahead of logistic regression across
the evaluated training sizes, and the Lending Club random-split ordering is
reproduced under temporal evaluation.


\section{Relation to Prior Findings}

\citet{bahloolPerformanceFairnessExplainability2026} characterise the
performance--explainability trade-off in credit scoring as more often assumed than
empirically quantified, because performance and explainability are commonly studied
separately. The present findings support that diagnosis. An ensemble advantage over the EBM resolves only on Lending Club, whereas
one ROC-AUC comparison on Home Credit resolves in the opposite direction and the
remaining comparisons show no consistent ensemble advantage. This establishes no equivalence. It does shift the burden of argument: a claim that
intrinsic interpretability entails a predictive cost requires evidence in the setting
where it is made and comparison with functionally flexible glass-box alternatives
rather than linear baselines alone.

The observed magnitudes are broadly consistent with the closest quantitative reference
points. Across 45 tabular datasets,
\citet{amekoeExploringAccuracyInterpretability2024b} report an average relative
accuracy loss below $4\%$ for intelligible models.
\citet{dessainCostExplainabilityAI2023} report fifteen to twenty basis points of
return between their best black-box and interpretable models, against about five
basis points of normalised principal here. These quantities are not directly
comparable. The first measures accuracy across a wider tabular benchmark, while
the second measures realised return rather than EMP. They nevertheless support the
same broader interpretation: where a performance cost exists, its magnitude must be
evaluated in context rather than inferred from the model family alone.

The Lending Club results contextualise the findings of \citet{schwartzEnhancingMLModels2025}, 
who report the EBM narrowly outperforming an XGBoost baseline on the same raw dataset 
(ROC-AUC $0.674$ versus $0.669$). The present study instead yields a consistent but small 
ensemble advantage in both ROC-AUC and economic value. This contrast does not necessarily 
indicate contradictory model behaviour. Rather, it highlights how sensitive relative 
performance is to specific pipeline configurations and empirical contexts. Whether the 
EBM slightly leads or slightly trails the tree ensembles ultimately depends on design 
choices such as feature selection and evaluation criteria. The shared conclusion across 
both pipelines, however, is that the EBM performs at a highly competitive level. It 
therefore emerges not as a mere transparency-oriented fallback, but as a primary model 
candidate whose exact rank must be established in the target context.

\section{Regulatory Context}\label{sec:regulatory-reading}

The framework set out in Section~\ref{sec:regulatory} does not translate the
applicable legal requirements into a preference for a particular model family.
Neither the AI Act nor the GDPR makes an EBM compliant by construction or a
boosted ensemble non-compliant because its explanations are generated post hoc.
The legal assessment instead concerns the deployed system, including its intended
use, documentation, data governance, human oversight, monitoring and the
information provided to affected persons. Model architecture may supply relevant
technical evidence but cannot replace this system-level assessment.

The evaluated model families nevertheless differ structurally in how explanation
evidence is generated. Logistic regression and the EBM expose components of their
fitted prediction functions, whereas XGBoost and CatBoost require a separate
TreeSHAP computation. As \citet{liptonMythosModelInterpretability2018} stresses,
a post-hoc explanation is not equivalent to direct transparency of the fitted
model. This distinction affects how model behaviour can be documented and
examined, but it alone establishes neither legal sufficiency nor a ranking of
compliance effort. Direct inspectability does not prove compliance, while reliance
on TreeSHAP does not prove non-compliance.

The empirical metrics reported in this thesis may inform the assessment of a
deployed system, but none constitutes a statutory threshold or conformity finding.
Accordingly, neither architecture is excluded on model-family grounds, and this
thesis does not assess whether any evaluated pipeline satisfies the applicable
requirements.

\section{Implications for Model Choice}

The central implication is procedural: model selection should not begin from a
presumed performance penalty for intrinsic interpretability. While 
\citet{amekoeExploringAccuracyInterpretability2024b} recommend incorporating inherently 
interpretable models primarily when transparency is a key constraint, the marginality of 
the trade-offs demonstrated in this thesis suggests that the EBM should enter the 
initial candidate set alongside boosted-tree ensembles from the outset. Rather than 
serving as a transparency-oriented fallback after a black-box model has already 
been selected, the glass-box model must be evaluated on its own predictive and 
economic merits. Its direct inspectability is an additional structural property, 
not its sole justification.


\section{Limitations}\label{sec:limitations}

Several design choices constrain the generalisability of the findings and define
the limits of the evaluated empirical setup.

\paragraph{Data and population.}
All four datasets contain only \emph{booked} loans. The findings are therefore
conditional on historical acceptance decisions and do not evaluate model performance
for the through-the-door applicant population. No reject-inference procedure was
applied. Moreover, the datasets are historical public benchmarks and represent
contemporary operational credit scoring only partially. Changes in lending products,
underwriting policies, macroeconomic conditions, applicant behaviour and the
availability of decision-time information may alter both predictive relationships and
relative model performance. The findings should therefore be interpreted as evidence
under the evaluated benchmark conditions rather than as expected performance for a
current institution-specific applicant population. No group-wise fairness evaluation was conducted, so the results support no conclusions
about disparate outcomes across applicant groups.

\paragraph{Model coverage.}
Each side of the comparison rests on limited model families. The flexible side is
represented by gradient-boosted trees in two implementations, chosen because they are
the documented reference for tabular credit data, but neural architectures, random
forests and other flexible learners were not evaluated. The intrinsically interpretable
side is represented by one non-linear model, so the within-family check available for
the ensembles has no counterpart there. Findings on the EBM consequently may not transfer
to other interpretable classes and
the conclusions concern the four evaluated pipelines rather than the model classes as a
whole.

\paragraph{Feature representation.}
Feature encoding is not neutral between model families. Ordinal encodings represent
ordered variables in a form that the models may use differently, while aggregating
one-hot encoded columns into source-feature importances can affect the resulting
attribution rankings. These effects were not isolated experimentally. The reported
comparisons therefore partly reflect interactions between the algorithms and the
chosen feature representation rather than model-class differences alone.

\paragraph{Model selection and tuning.}
All configurations were selected to maximise validation ROC-AUC. The EMP results
therefore describe the economic performance of pipelines selected for their overall
discriminatory ability rather than pipelines optimised directly for expected profit.
These objectives need not select the same configuration: ROC-AUC evaluates ranking
across the full range of decision thresholds, whereas EMP rewards ranking performance
in the regions that are economically relevant under its assumed returns, losses and
class prevalence. A configuration with a slightly lower ROC-AUC could consequently
produce a higher EMP. Direct EMP-based tuning might therefore change
both the selected configurations and the relative ordering of the model families. The
profit-based comparison is conditional on the discrimination-first selection strategy
and does not establish which models would perform best under direct economic
optimisation.

\paragraph{Uncertainty and temporal coverage.}
The headline results rest on one train--test split per dataset. The paired bootstrap
quantifies test-set sampling uncertainty conditional on the fitted pipelines. It does
not propagate uncertainty from alternative splits, repeated model selection or
alternative explainer reference samples. No uncertainty intervals were calculated for
between-pipeline differences in attribution-ranking stability, so their relative
ordering remains descriptive.

Out-of-time evaluation was possible only for Lending Club and uses a single
chronological design rather than repeated rolling-origin assessments. It is also
subject to maturity-dependent censoring because later cohorts contain a smaller
proportion of loans with resolved outcomes. The temporal analysis therefore cannot
separate changes caused via temporal drift from those caused by changing outcome
availability, nor does it establish that the observed ordering persists across
successive origination cohorts or under changing market conditions.

\paragraph{Metrics and explanations.}
EMP remains conditional on the assumed return, loss-given-default distribution and
class prevalence. The use of a common loss distribution and partly standardised return
assumptions enables a consistent experimental comparison across the datasets, but it
does not reproduce the specific economics of each credit product or lending
institution. Product-specific estimates of returns and loss severity could shift the
economically relevant operating region and thereby change not only absolute EMP levels,
but also the pairwise differences and potentially the model ordering. The reported EMP
results therefore describe economic performance under the stated assumptions rather
than realised portfolio profit or a definitive profitability ranking
\citep{xuProfitRiskdrivenCredit2024}.

The explanation analysis combines model-specific methods with different reference
constructions and decomposition rules. Its results consequently describe
model--explainer pipelines rather than model properties alone. Attribution-ranking
stability measures the reproducibility of global rankings under case resampling. It
does not test whether those rankings faithfully represent the features governing the
fitted model's predictions. Human intelligibility and the usefulness of local
explanations to credit decision-makers or applicants were likewise not evaluated.

\paragraph{Evidence base.}
The closest prior comparison, \citet{schwartzEnhancingMLModels2025}, is a preprint. It
was retained because it is the only identified study evaluating the EBM against
XGBoost on Lending Club, and omitting the closest precedent would have left the
comparison without a direct reference point. Its reported values should nevertheless be
treated as unverified. Two circumstances limit the resulting reliance. First, the
preprint enters the discussion as a contrast rather than as support: it reports the EBM
narrowly ahead of XGBoost, whereas the present study finds a small ensemble advantage on
the same dataset. Second, the broader interpretation that flexible interpretable models
can be competitive is carried by peer-reviewed evidence
\citep{amekoeExploringAccuracyInterpretability2024b}.


\section{Future Work}\label{sec:future}

A deployment-oriented extension should evaluate the pipelines on current institutional
credit data collected at a consistent decision point. Reject-inference and sensitivity
analyses could examine whether the relative model performance extends beyond the booked
population, while prespecified group-fairness
analyses could test whether the observed performance differences persist across
relevant applicant groups. Product-specific
estimates of returns, loss severity and class prevalence would additionally indicate
whether the economic ordering obtained under the standardised EMP assumptions transfers
to operational portfolios.

Repeating the full pipeline across independently drawn data partitions would propagate
sources of uncertainty held fixed in the present design and permit direct estimation of
between-pipeline differences in attribution-ranking stability.

EMP-based or
multi-objective hyperparameter selection could test whether discrimination-first tuning
conceals configurations with preferable economic trade-offs. Temporal validation also
warrants extension to additional datasets with reliable origination timestamps and
sufficiently mature outcomes.

Extending the model set would test how broadly the observed performance pattern applies.
Additional glass-box approaches could show whether the EBM's competitiveness is specific
to its architecture or extends to other intrinsically interpretable models. Including
flexible predictors beyond boosted trees would likewise test whether the findings hold
against a larger set of black-box alternatives.

Explanation evaluation should distinguish stability from faithfulness and practical
usefulness. Retraining- or perturbation-based tests designed to limit unrealistic
out-of-distribution inputs could assess whether the attributed features accurately
represent the fitted model's predictive behaviour
\citep{zhengFFidelityRobustFramework2025}. Local explanations could then be evaluated
with credit decision-makers and affected applicants to determine whether they are
understandable, decision-relevant and useful for understanding or contesting individual
outcomes \citep{nautaAnecdotalEvidenceQuantitative2023}.

Finally, a controlled organisational study could examine whether direct inspectability
changes the resources required to document, validate, maintain and communicate a
credit-scoring system relative to a separate post-hoc explanation pipeline. Relevant
outcomes include development and validation time, documentation obligations,
explainer-specific failure modes and stakeholder comprehension. The structural and
regulatory differences between inherent and post-hoc explanations motivate this
comparison but do not predetermine its result.

